\section{A Loser-Side Screen for Cartel Adjacency}
\label{sec:award_screen}

\subsection{The Participation Primitive}
\label{sec:award_primitive}

The award layer cannot reveal why a firm lost. A losing bid may
be sincere, poorly priced, capacity constrained, exploratory,
strategic, or coordinated; the analyst querying award records
cannot distinguish these from one tender's outcome. What the
award layer can observe is thinner: which firms appeared, which
firms won, and how often the same firm kept appearing without
ever winning. The primitive used here is therefore not a
behavioral type. It is an observable loser-side exposure measure
built from participation and zero wins. In the BEC data the
relevant conditioning set contains $\valAlwaysLosers$
always-loser firms.

Zero wins alone is not enough. Many firms never win for
ordinary competitive reasons. They may be too small or poorly
matched to a particular item; they may face binding capacity
constraints in the periods when they would otherwise have won;
they may be using procurement participation for market
intelligence, eligibility maintenance, subcontracting
positioning, or relationship-building with procuring units.
Without an additional discriminating signal, the always-loser
stratum is too heterogeneous for triage. The screening signal is
persistence inside that stratum: among firms that never win,
how often do they keep showing up?

The maintained condition is monotonic rather than
classificatory: within the zero-win stratum, cartel-deployed
losing participants, if present, should participate more
persistently than ordinary losing firms, because their function
in the cartel---supplying the appearance of competition while
leaving the designated winner in place---pays off through
repeated deployment in cartel-targeted tenders. Under that
condition, participation intensity inside the zero-win stratum,
implemented as $\log(1+\text{tenders\_count})$, is a monotone
ranking statistic for loser-side cartel-adjacency risk. The
appendix formalizes this condition in a sparse screening
framework; the main text uses it only to discipline the ranking
interpretation.

This primitive does not identify cover bidders or cartel
members. It does not establish agreement, liability, or damages.
It defines a first-stage ranking: among firms the award layer
only observes as losers, which ones appear persistently enough
to justify bid-layer forensic attention? The binary rule in
subsection~\ref{sec:award_rule} is only an operational
coarsening of this continuous ranking.

\subsection{The Binary Rule}
\label{sec:award_rule}

The continuous ranking is the economic object. Within the
always-loser stratum, we score firms by
$\log(1+\text{tenders\_count})$, which orders repeated zero-win
participation under the monotonicity condition stated in
subsection~\ref{sec:award_primitive}.

A binary frequent-loser flag is useful for a different reason.
Enforcement programs need more than a rank: agencies must
decide which firms enter a survivor pool, which records justify
a costly bid-microdata recovery, and where a forensic queue
begins on Monday morning. The cutoff converts the continuous
score into an administrable first-stage list, with the
auditable property that two analysts applying the same rule to
the same participation distribution will obtain the same
survivor pool.

For the discrete rule, we use the median plus $1.5\times$
interquartile range of the always-loser participation
distribution. In BEC this gives a cutoff of $\valThreshold$
tender-items with zero wins and flags $\valFL$ frequent losers.
The rule is fixed from the participation distribution before
any CADE label enters the analysis and requires no bid amounts,
bid ranks, or within-tender dispersion measures. The cutoff
therefore inherits the low-information property that motivates
the architecture.

The number $\valThreshold$ has no independent legal or
structural status. It does not divide cartel firms from
non-cartel firms, and it does not make participants below the
threshold competitive or those above it liable. It is an
auditable implementation of a continuous ranking, calibrated to
the spread of zero-win participation in the data rather than to
any structural parameter of the bidding environment.

Agencies could implement the same architecture differently. A
capacity-constrained top-$k$ list would draw a fixed number of
firms from the highest-ranked end of the continuous score; a
budget-calibrated threshold would set the cutoff to match the
bid-microdata recovery capacity of an investigative team; a
dynamically refreshed rule would update the threshold as new
participation records arrive. The paper reports both the binary
flag and the continuous score where the analysis allows, so
that the validation in Section~\ref{sec:validation} can be read
as evidence on the ranking architecture rather than on any one
cutoff. Threshold manipulability by firms aware of the rule is
a practical concern about the binary implementation, not a
conceptual collapse of the architecture, and we return to it
under adversarial adaptation.

\subsection{Populations and Legal-Economic Role}
\label{sec:award_populations}

The empirical design turns on keeping five populations
separate. The broad BEC universe defines where procurement
participation is observed. The always-loser stratum defines the
side of the bidding pool that the award layer can rank.
Frequent losers are the administrable subset flagged by the
screen. Direct CADE defendants are legal membership anchors from
adjudicated cases. Cobidders are always-loser firms observed
alongside those anchors. These populations overlap in the data,
but they do not play the same legal or empirical role.

Table~\ref{tab:populations_v18} fixes those roles. The screen
operates only inside the always-loser stratum. Frequent-loser
status is an operational flag, not a legal category. Direct
defendants locate cartel-affected environments, but a zero-win
loser-side statistic is not designed to classify them. Cobidders
define the cartel-adjacency target for validation: they are
firms that lose, keep appearing, and are observed near
adjudicated cartel defendants. That label is informative for
validation but is not proof of agreement, membership, or
liability.

\begin{table}[htbp]
\centering
\caption{Populations and legal-economic role. Frequent losers
are an operational screen; direct defendants are legal anchors;
cobidders are cartel-adjacent loser-side labels. The table
separates legal membership from screening exposure.}
\label{tab:populations_v18}
\footnotesize
\begin{adjustbox}{max width=\textwidth,center}
\begin{tabular}{p{4.0cm}p{2.1cm}p{7.0cm}}
\toprule
Population & Size & Role in the analysis \\
\midrule
All BEC participants & $41{,}444$ & Firms participating in BEC tenders during $2009$--$2019$; the outer procurement universe. \\
Always-losers & $\valAlwaysLosers$ & Firms with zero wins; the loser-side stratum within which the award-layer screen ranks firms. \\
Frequent losers & $\valFL$ & Always-losers above the participation cutoff; the discrete first-stage screen. \\
CADE direct defendants in BEC & $\valDirectCADE$ & Firms named in CADE procurement-cartel rulings and active in BEC; the legal membership anchors. \\
Cobidders in the always-loser stratum & $\valCobidders$ & Always-losers that participated alongside direct defendants; the cartel-adjacency target for validation. \\
\bottomrule
\end{tabular}
\end{adjustbox}
\end{table}

This separation is what makes the validation design legally
limited rather than opportunistic. The paper does not collapse
defendants and cobidders into one cartel class. Frequent-loser
status is not cartel membership; cobidder status is not legal
membership; direct-defendant status is the legal-anchor role,
not the screen target. The screen ranks loser-side exposure;
CADE labels anchor the validation step.

Section~\ref{sec:validation} therefore validates the ranking,
not the legal status of the ranked firms. It asks whether the
frequent-loser ranking concentrates on the cartel-adjacent
loser-side population at rates above what participation volume
alone would imply.
